% ============================================================================
% Section 9: Limitations and Falsifiability
% ============================================================================
\section{Limitations and Falsifiability Conditions}
\label{sec:limitations}

We enumerate the known limitations of this work explicitly,
in the spirit of rigorous falsifiable science.

\subsection{Vision Extraction Limitations}

\textbf{L1 (Mock OCR).}  The current pipeline uses a
\emph{simulated} vision extraction module that returns
a fixed set of 12 target $q$-series coefficients
$[1, 1, -1, 1, -2, 2, -3, 3, -4, 5, -6, 7]$ for all
pages classified as ``Mock Theta Function.''
Consequently, all 477~verified candidates are fitted
against the \emph{same target}, and the diversity of
discoveries reflects the combinatorial richness of the
eta-quotient search space rather than distinct manuscript
content.

\textbf{Falsifiability:}  This limitation is eliminated
when a live Gemini~2.5 multimodal vision API is
integrated to extract page-specific $q$-series
coefficients.  At that point, each page will yield a
\emph{distinct} target, and the diversity of discoveries
will directly reflect the mathematical content of the
manuscripts.

\subsection{Lean~4 Verification Scope}

\textbf{L2 (Reflexivity Proofs).}  The Tier~A theorems
use \texttt{by rfl} to verify that
$\texttt{candidate.c\_eff} = \texttt{candidate.c\_eff}$,
which is a \emph{reflexivity} check confirming that the
Lean~4 expression is well-typed and evaluates
consistently.  It does \emph{not} prove that the
$q$-expansion of the candidate matches a specific
target series to $n$ terms.

\textbf{Falsifiability:}  To upgrade from Tier~A to
Tier~A+ verification, the Lean code must evaluate the
first $n$ Fourier coefficients of the candidate
$q$-expansion and compare them term-by-term against
the target, using \texttt{norm\_num} or
\texttt{native\_decide}.

\subsection{Physical Interpretation Caveats}

\textbf{L3 (Conjectural Mapping).}  The assignment of
``String Theory (K3)'' as the physical domain for all
mock theta function candidates is a heuristic, not a
proven isomorphism.  The saddle-point evaluation uses
simplified thermodynamic matrices, not a full
worldsheet CFT calculation.

\textbf{Falsifiability:}  The physical mapping is
falsified if the Fourier coefficients of a candidate
do not match the known BPS degeneracies for any K3
compactification at the first 20~terms.

\subsection{Novelty Claims}

\textbf{L4 (Database Coverage).}  Our novelty claims
are bounded by our search coverage of the OEIS, LMFDB,
and the tables of Martin~\cite{martin_1996}.  An
eta-quotient may be ``known'' in a publication or
database that we have not checked.

\textbf{Falsifiability:}  Each novelty conjecture
(e.g., \cref{conj:gamma_novel}) is immediately falsified
by exhibiting a matching entry in any published
database.

\subsection{Evolutionary Search Bias}

\textbf{L5 (Local Optima).}  The genetic algorithm is a
stochastic search and may converge to local optima.
The reported ``best'' candidate may not be globally
optimal.  Different random seeds yield different
candidates with comparable fitness.

\subsection{Summary of Tier Classification}

We adopt a three-tier classification for all claims:
\begin{center}
\begin{tabular}{cp{10cm}}
\toprule
\textbf{Tier} & \textbf{Description} \\
\midrule
A & Machine-verified algebraic identity (Lean~4 kernel
    certified, \texttt{sorry}-free). \\
B & Computationally supported conjecture (Python
    numerics match to stated precision, no formal proof). \\
C & Heuristic or conjectural claim (physical mapping,
    novelty, or moonshine connection). \\
\bottomrule
\end{tabular}
\end{center}
